\chapter{El grafo y el ciclo hamiltoniano}
\label{cap:grafo}

\section{Definiciones básicas}
Un \textbf{grafo} $G=(V,E)$ es un conjunto de vértices $V=\{1,\dots,n\}$ unidos por aristas $E$.
En este proyecto los grafos son \emph{no dirigidos} (la arista $\{u,v\}$ sirve en los dos
sentidos) y \emph{simples} (sin lazos ni aristas repetidas). El \textbf{grado} de un vértice es el
número de aristas que llegan a él.

Un \textbf{camino} es una secuencia de vértices en la que cada par consecutivo está unido por una
arista. Un \textbf{ciclo} es un camino que termina donde empezó. Un \textbf{ciclo hamiltoniano}
es un ciclo que pasa por \emph{todos} los vértices exactamente una vez. Dicho de otra forma, es
una lista $(v_1,v_2,\dots,v_n)$ con estas dos condiciones:
\begin{enumerate}
  \item Contiene cada vértice del grafo una sola vez, es decir, es una \emph{permutación} de $V$.
  \item Cada vértice está conectado con el siguiente, y el último con el primero:
        $\{v_i,v_{i+1}\}\in E$ para $i=1,\dots,n-1$, y además $\{v_n,v_1\}\in E$.
\end{enumerate}
Por eso un ciclo hamiltoniano en un grafo de $n$ vértices usa exactamente $n$ aristas: las $n-1$
del camino más la arista que lo cierra.

\section{Por qué es un problema difícil}
Decidir si un grafo tiene ciclo hamiltoniano es uno de los problemas NP-completos clásicos de
Karp \cite{de1}. En la práctica, esto significa que no conocemos un algoritmo rápido para
todos los casos. Además, el número de rutas posibles crece como $(n-1)!/2$: con 20 vértices ya son
más de $6\times10^{16}$, así que revisarlas una por una no es opción.

Es importante distinguir dos cosas. Si un algoritmo heurístico \emph{encuentra} un ciclo, podemos
comprobarlo fácilmente y la respuesta es segura. Pero si \emph{no} lo encuentra, eso no prueba
que el grafo no tenga ciclo: solo significa que se acabó el tiempo. Por eso, en este reporte,
una corrida sin éxito la llamamos \emph{inconclusa}.

El problema está muy relacionado con el TSP \cite{sa3}: si a las aristas del grafo les damos peso
0 y a las que faltan peso 1, un recorrido del TSP con costo 0 es justamente un ciclo hamiltoniano.
Esta idea es la que usamos para calificar las rutas (sección~\ref{cap:repr}).

\section{Ejemplo}
Tomemos un grafo con 5 vértices y las aristas
$\{1,2\},\{2,3\},\{3,4\},\{4,5\},\{5,1\},\{1,3\}$.

\begin{itemize}
  \item La ruta $(1,2,3,4,5)$ \textbf{sí} es un ciclo hamiltoniano: $1\!-\!2$, $2\!-\!3$,
        $3\!-\!4$ y $4\!-\!5$ existen, y la arista de cierre $5\!-\!1$ también.
  \item La ruta $(1,3,2,4,5)$ \textbf{no} lo es, aunque visita todos los vértices: no existe la
        arista $2\!-\!4$.
\end{itemize}

\section{Cómo guardamos el grafo en el programa}
Los grafos del FHCP vienen en el formato HCP, que es parte del estándar TSPLIB. El archivo
indica el número de vértices (\texttt{DIMENSION}) y luego la lista de aristas
(\texttt{EDGE\_DATA\_SECTION}):

\begin{verbatim}
NAME : graph4.hcp
COMMENT : FHCP Challenge Set. http://fhcp.edu.au/fhcpcs
TYPE : HCP
DIMENSION : 84
EDGE_DATA_FORMAT : EDGE_LIST
EDGE_DATA_SECTION
1 43
1 46
1 58
...
76 77
-1
EOF
\end{verbatim}

El archivo \texttt{src/grafo.py} lee ese formato y guarda el grafo como una \emph{lista de
adyacencia}: para cada vértice, un conjunto con sus vecinos. Así, saber si dos vértices están
conectados es una consulta inmediata. Internamente numeramos los vértices de $0$ a $n-1$ y los
pasamos de vuelta a $1..n$ al mostrar o guardar los resultados.

Cuando un algoritmo encuentra un ciclo, lo guardamos en formato TOUR, que es el mismo formato en
que el FHCP publica sus soluciones oficiales (\texttt{TOUR\_SECTION} seguida de la lista de
vértices y un $-1$ al final). Esto nos permitió comparar directamente nuestras soluciones con las
oficiales.
